\documentclass[../main.tex]{subfiles}

\graphicspath{{../image/}{../figure/}}

\begin{document}

\chapter{生成对抗网络}\label{chapter:ch17-生成对抗网络}

生成式模型使用机器学习算法从一组训练数据中学习潜在的数据分布，然后根据该分布生成新的样本。例如，一个生成式模型可以基于动物图像进行训练，随后用于生成新的动物图像。我们可以将这种生成式模型理解为一个条件分布 $p(x\mid w)$，这里 $x$ 是数据空间中的向量，而 $w$ 代表模型的可学习参数。在很多情况下，我们更倾向于使用形式为 $p(x\mid c,w)$ 的条件生成式模型，其中 $c$ 代表条件变量的向量。以动物图像生成模型为例，我们希望生成的图像应该是特定动物的，比如猫或狗，而这是由 $c$ 的值指定的。

对于像图像生成这样的实际应用而言，其分布极其复杂，深度学习能显著提高生成式模型的性能。在讨论基于 Transformer 的自回归大语言模型时，我们已经了解了一些重要的深度生成式模型。我们还概述了 4 种基于非线性潜变量模型的生成式模型（参见\cref{chapter:ch12-Transformer与注意力机制}和\cref{subsection:ch16-四种生成模型}），本章将讨论其中的第一种——生成对抗网络（Generative Adversarial Network, GAN），其他三种将在后续各章中讨论。

\subfile{Chapter-17/section-01.tex}
\subfile{Chapter-17/section-02.tex}
\subfile{Chapter-17/section-03.tex}
\subfile{Chapter-17/section-04.tex}
\subfile{Chapter-17/section-05.tex}

\end{document}
